\chapter[IA et Gestion de la Qualité]{Intelligence Artificielle et Gestion de la Qualité}
\label{chap:ia_quality}
\chapterplan{%
\chapterplanitem{}{Introduction}{sec:chap5_introduction}
\chapterplanitem{1}{Sprint 3 : Moteur d'IA et scoring des appels}{sec:sprint3}
\chapterplanitem{2}{Sprint 4 : Gestion de la qualité et évaluations}{sec:sprint4}
\chapterplanitem{}{Conclusion}{sec:chap5_conclusion}
}

\section*{Introduction}
\phantomsection
\label{sec:chap5_introduction}

Ce chapitre présente la deuxième étape fonctionnelle majeure du projet. Après la mise en place
du socle d'authentification et de gestion des leads, le travail s'oriente vers le coeur
innovant de la plateforme~: l'intelligence artificielle pour l'analyse et le scoring des
appels, ainsi que la gestion de la qualité.

Le Sprint 3 couvre l'intégration du moteur d'IA~: transcription automatique des appels via
Whisper, scoring IA basé sur gemma3:4b avec des critères pondérés, et analyse de sentiment
via DistilBERT. Le Sprint 4 complète ce périmètre avec l'évaluation manuelle de la qualité,
les alertes et les tableaux de bord de performance.

%% ══════════════════════════════════════════════════════════════════════════════
\section[Sprint 3 : Moteur d'IA]{Sprint 3 : Moteur d'IA et scoring des appels}
\label{sec:sprint3}

%% ────────────────────────────────────────────────────────────────────────────
\subsection{Objectifs et Backlog du Sprint 3}
\label{subsec:sprint3_backlog}

Le Sprint 3 a pour objectif de livrer le moteur d'intelligence artificielle de la plateforme.
Il permet la transcription automatique des appels audio, le scoring qualitatif basé sur des
critères pondérés et l'analyse de sentiment.

\begingroup
\small
\setlength{\tabcolsep}{5pt}
\renewcommand{\arraystretch}{1.35}
\begin{longtable}{|>{\centering\arraybackslash}p{1.1cm}|>{\raggedright\arraybackslash}p{9.0cm}|>{\centering\arraybackslash}p{1.6cm}|>{\centering\arraybackslash}p{1.6cm}|}
\caption{Backlog du Sprint 3 -- Moteur d'IA et scoring}\label{tab:sprint3_backlog}\\
\hline
\textbf{ID} & \textbf{User Story} & \textbf{Priorité} & \textbf{Statut} \\
\hline
\endfirsthead
\hline
\textbf{ID} & \textbf{User Story} & \textbf{Priorité} & \textbf{Statut} \\
\hline
\endhead
\hline
\endfoot
\endlastfoot
4.1 & En tant que \textbf{système}, je veux transcrire automatiquement les appels via Whisper afin de produire un texte exploitable. & Élevée & Terminé \\
\hline
4.2 & En tant que \textbf{système}, je veux scorer automatiquement chaque appel via l'IA afin d'évaluer sa qualité. & Élevée & Terminé \\
\hline
4.3 & En tant qu'\textbf{Agent}, je veux consulter mon score IA par appel afin de connaître ma performance. & Élevée & Terminé \\
\hline
4.4 & En tant que \textbf{Quality Manager}, je veux visualiser les scores IA détaillés par critère afin d'analyser la qualité. & Élevée & Terminé \\
\hline
4.5 & En tant que \textbf{Super Admin}, je veux configurer les poids des critères de scoring afin d'adapter l'évaluation. & Élevée & Terminé \\
\hline
4.6 & En tant que \textbf{système}, je veux analyser le sentiment des appels via DistilBERT afin de détecter les tendances. & Moyenne & Terminé \\
\hline
4.7 & En tant que \textbf{Supervisor}, je veux consulter un dashboard IA agrégé afin de visualiser les performances globales. & Élevée & Terminé \\
\hline
4.8 & En tant que \textbf{Quality Manager}, je veux comparer le score IA au score manuel afin de valider la pertinence. & Moyenne & Terminé \\
\hline
\end{longtable}
\endgroup

%% ────────────────────────────────────────────────────────────────────────────
\subsection{Spécification et Conception du Sprint 3}
\label{subsec:sprint3_spec}

\subsubsection{Architecture du pipeline IA}

Le pipeline IA est composé de trois étapes principales~:

\begin{enumerate}
    \item \textbf{Transcription (Whisper)}~: le fichier audio de l'appel est transmis au
    modèle Whisper, qui produit une transcription textuelle en français.
    \item \textbf{Scoring (gemma3:4b)}~: la transcription est envoyée au LLM avec un prompt
    structuré demandant l'évaluation de 8 critères pondérés (accueil, énergie, voix, écoute,
    client, opérateur, efficacité, conclusion). Le LLM retourne un score composite sur 100.
    \item \textbf{Analyse de sentiment (DistilBERT)}~: la transcription est analysée pour
    déterminer le sentiment global (positif, négatif, neutre).
\end{enumerate}

\subsubsection{Configuration du scoring}

Les poids des critères de scoring sont configurables via un fichier \texttt{config.json}~:

\begin{center}
\begin{tabular}{|l|c|}
\hline
\textbf{Critère} & \textbf{Poids (\%)} \\
\hline
Accueil & 10 \\
\hline
Énergie & 15 \\
\hline
Voix & 10 \\
\hline
Écoute & 15 \\
\hline
Client & 15 \\
\hline
Opérateur & 15 \\
\hline
Efficacité & 10 \\
\hline
Conclusion & 10 \\
\hline
\end{tabular}
\end{center}

\subsubsection{Diagramme de séquence -- Scoring d'un appel}

\begin{figure}[H]
    \centering
    \fbox{\parbox{0.95\linewidth}{
    \centering\vspace{0.4cm}
    \textit{[Diagramme de séquence : Pipeline de scoring IA]}\\[0.3cm]
    L'agent termine un appel, le fichier audio est sauvegardé, le service AiScoringService
    déclenche la transcription via Whisper, puis envoie la transcription au LLM gemma3:4b
    via Ollama avec les poids configurés, reçoit le score et persiste le résultat dans
    la base de données. L'agent peut ensuite consulter son score depuis le dashboard.
    \vspace{0.4cm}
    }}
    \caption{Diagramme de séquence -- Pipeline de scoring IA}
    \label{fig:sprint3_ia_pipeline}
\end{figure}

%% ────────────────────────────────────────────────────────────────────────────
\subsection{Présentation des Technologies Utilisées}
\label{subsec:sprint3_tech}

\subsubsection{Whisper pour la transcription audio}

Whisper est un modèle de reconnaissance vocale développé par OpenAI~\cite{Whisper}. Il est
utilisé avec la configuration suivante~:

\begin{itemize}
    \item \textbf{Modèle}~: \texttt{base} (optimisé pour le français)
    \item \textbf{Langue}~: français
    \item \textbf{Exécution}~: locale (pas d'appel API externe)
\end{itemize}

\subsubsection{Ollama avec Gemma 3 pour le scoring}

Ollama est la plateforme d'exécution locale de LLM~\cite{Ollama}. Le modèle gemma3:4b de
Google DeepMind~\cite{Gemma3} est utilisé avec les paramètres suivants~:

\begin{itemize}
    \item \textbf{Température}~: 0.3 (faible créativité pour des évaluations cohérentes)
    \item \textbf{Timeout}~: 120 secondes
    \item \textbf{Prompt structuré}~: un prompt détaillé demande l'évaluation de chaque
    critère avec justification et retourne un score JSON structuré.
\end{itemize}

Le prompt système utilisé pour le scoring est conçu pour obtenir une évaluation standardisée~:

\begin{quote}
\small\textit{``Tu es un évaluateur de qualité pour centre d'appels. Analyse la transcription
suivante et note chaque critère de 0 à 100 selon les poids définis. Retourne UNIQUEMENT
un objet JSON avec les scores détaillés et le score total.''}
\end{quote}

\subsubsection{DistilBERT pour l'analyse de sentiment}

DistilBERT~\cite{DistilBERT} est utilisé pour classifier le sentiment des transcriptions en
trois catégories~: positif, négatif ou neutre.

\subsubsection{Le service de scoring IA}

Le service \texttt{AiScoringService} orchestre le pipeline complet~:

\begin{itemize}
    \item \textbf{Scoring}~: envoie la transcription au LLM avec les poids configurés et
    parse la réponse JSON pour extraire les scores par critère.
    \item \textbf{Éligibilité}~: calcule le score d'éligibilité (\texttt{EligibilityCalculator})
    basé sur les règles métier.
    \item \textbf{Historique}~: conserve un log de chaque scoring dans
    \texttt{AiEligibilityLog} pour la traçabilité.
\end{itemize}

%% ────────────────────────────────────────────────────────────────────────────
\subsection{Réalisation}
\label{subsec:sprint3_realisation}

\subsubsection{Interface de scoring IA}

L'interface de scoring IA permet à l'agent de consulter ses scores par appel, avec le détail
par critère et l'évolution dans le temps.

\begin{figure}[H]
    \centering
    \fbox{\parbox{0.9\linewidth}{
    \centering\vspace{0.4cm}
    \textit{[Capture d'écran : Scores IA par appel]}\\[0.3cm]
    Cette capture présente les scores IA détaillés pour un appel~: score total sur 100,
    scores par critère (accueil, énergie, voix, écoute, client, opérateur, efficacité,
    conclusion), la transcription de l'appel et l'analyse de sentiment.
    \vspace{0.4cm}
    }}
    \caption{Scores IA par appel}
    \label{fig:sprint3_ia_scores}
\end{figure}

\begin{figure}[H]
    \centering
    \fbox{\parbox{0.9\linewidth}{
    \centering\vspace{0.4cm}
    \textit{[Capture d'écran : Dashboard IA agrégé]}\\[0.3cm]
    Le dashboard IA présente les métriques globales~: score moyen de l'équipe, évolution
    sur la période, répartition des scores, nombre d'appels analysés et tendances par
    critère.
    \vspace{0.4cm}
    }}
    \caption{Dashboard IA agrégé}
    \label{fig:sprint3_ia_dashboard}
\end{figure}

\subsubsection{Configuration des poids de scoring}

L'interface de configuration permet au Super Admin d'ajuster les poids de chaque critère
en temps réel, avec validation que la somme des poids est égale à 100.

\begin{figure}[H]
    \centering
    \fbox{\parbox{0.9\linewidth}{
    \centering\vspace{0.4cm}
    \textit{[Capture d'écran : Configuration des poids de scoring]}\\[0.3cm]
    Interface de configuration des poids avec des curseurs pour chaque critère
    (accueil, énergie, voix, écoute, client, opérateur, efficacité, conclusion).
    Un indicateur visuel confirme que la somme des poids est correcte.
    \vspace{0.4cm}
    }}
    \caption{Configuration des poids de scoring IA}
    \label{fig:sprint3_weight_config}
\end{figure}

\subsubsection{Tests unitaires et validation}

Les tests du Sprint 3 couvrent~: le service de scoring IA (appel à Ollama, parsing des
réponses, gestion des timeouts), le service de transcription Whisper, le calculateur
d'éligibilité, les scénarios d'erreur (modèle indisponible, timeout, réponse invalide).

\subsubsection{Résultats obtenus}

Le Sprint 3 livre le moteur d'IA complet de la plateforme. La transcription Whisper permet
d'obtenir des textes exploitables des appels en français. Le scoring via gemma3:4b produit
des évaluations cohérentes et détaillées. Les agents disposent désormais d'un feedback
automatisé sur leur performance, et les managers d'une vue agrégée de la qualité des
interactions.

%% ══════════════════════════════════════════════════════════════════════════════
\section[Sprint 4 : Qualité]{Sprint 4 : Gestion de la qualité et évaluations}
\label{sec:sprint4}

%% ────────────────────────────────────────────────────────────────────────────
\subsection{Objectifs et Backlog du Sprint 4}
\label{subsec:sprint4_backlog}

Le Sprint 4 complète les fonctionnalités IA en permettant l'évaluation manuelle de la qualité
des appels, la configuration d'alertes et la visualisation de tableaux de bord de performance.

\begingroup
\small
\setlength{\tabcolsep}{5pt}
\renewcommand{\arraystretch}{1.35}
\begin{longtable}{|>{\centering\arraybackslash}p{1.1cm}|>{\raggedright\arraybackslash}p{9.0cm}|>{\centering\arraybackslash}p{1.6cm}|>{\centering\arraybackslash}p{1.6cm}|}
\caption{Backlog du Sprint 4 -- Gestion de la qualité}\label{tab:sprint4_backlog}\\
\hline
\textbf{ID} & \textbf{User Story} & \textbf{Priorité} & \textbf{Statut} \\
\hline
\endfirsthead
\hline
\textbf{ID} & \textbf{User Story} & \textbf{Priorité} & \textbf{Statut} \\
\hline
\endhead
\hline
\endfoot
\endlastfoot
5.1 & En tant que \textbf{Quality Manager}, je veux créer une évaluation manuelle d'un appel afin de noter la qualité. & Élevée & Terminé \\
\hline
5.2 & En tant que \textbf{Quality Manager}, je veux consulter l'historique des évaluations qualité afin de suivre les progrès. & Moyenne & Terminé \\
\hline
5.3 & En tant que \textbf{Super Admin}, je veux configurer des règles d'alerte qualité afin d'être notifié des dérives. & Élevée & Terminé \\
\hline
5.4 & En tant que \textbf{Quality Manager}, je veux recevoir une alerte quand un score passe sous un seuil. & Moyenne & Terminé \\
\hline
5.5 & En tant que \textbf{Supervisor}, je veux consulter un dashboard qualité afin de piloter la performance. & Élevée & Terminé \\
\hline
5.6 & En tant que \textbf{Quality Manager}, je veux exporter les résultats d'évaluation afin de les partager. & Moyenne & Terminé \\
\hline
\end{longtable}
\endgroup

%% ────────────────────────────────────────────────────────────────────────────
\subsection{Spécification et Conception du Sprint 4}
\label{subsec:sprint4_spec}

\subsubsection{Modèle d'évaluation qualité}

L'évaluation qualité repose sur deux sources complémentaires~:

\begin{itemize}
    \item \textbf{Scoring IA automatique}~: produit par le pipeline IA du Sprint 3 pour
    chaque appel. Score basé sur 8 critères pondérés.
    \item \textbf{Évaluation manuelle}~: réalisée par le Quality Manager. Elle comprend
    une note globale, des commentaires détaillés et peut inclure une note par critère.
\end{itemize}

Le score final composite est calculé en combinant les scores IA et manuels selon des règles
configurables, permettant d'obtenir une évaluation équilibrée et objective.

\subsubsection{Système d'alertes}

Le système d'alertes permet de configurer des seuils et des règles de déclenchement~:

\begin{itemize}
    \item \textbf{Score minimum}~: alerte déclenchée quand un score IA passe sous 70.
    \item \textbf{Inactivité}~: alerte après 15 minutes d'inactivité d'un agent.
    \item \textbf{Taux de conversion}~: alerte quand le taux de conversion passe sous 40\%.
\end{itemize}

Les alertes sont persistées dans la table \texttt{AlertHistory} et visibles dans le
dashboard du Quality Manager et du Supervisor.

\begin{figure}[H]
    \centering
    \fbox{\parbox{0.95\linewidth}{
    \centering\vspace{0.4cm}
    \textit{[Diagramme de séquence : Processus d'évaluation qualité]}\\[0.3cm]
    Le Quality Manager sélectionne un appel à évaluer, le système affiche la transcription
    et le score IA. Le Quality Manager saisit sa note et ses commentaires. Le service
    QualityService calcule le score composite et déclenche les vérifications d'alertes.
    \vspace{0.4cm}
    }}
    \caption{Diagramme de séquence -- Évaluation qualité}
    \label{fig:sprint4_quality_flow}
\end{figure}

%% ────────────────────────────────────────────────────────────────────────────
\subsection{Présentation des Technologies Utilisées}
\label{subsec:sprint4_tech}

\subsubsection{QualityScoreCalculator}

Le \texttt{QualityScoreCalculator} est un service dédié au calcul des scores composites.
Il prend en compte~:

\begin{itemize}
    \item Le score IA automatique avec ses sous-scores par critère.
    \item L'évaluation manuelle avec ses notes et commentaires.
    \item Les poids configurables pour chaque source d'évaluation.
\end{itemize}

\subsubsection{Système d'alertes et notifications}

Les alertes sont gérées par le service \texttt{InactivityAlertService} qui s'exécute en
arrière-plan. Il surveille périodiquement l'activité des agents et déclenche des alertes
selon les règles configurées. Les règles d'alerte sont définies dans la classe
\texttt{AlertRule} et supportent différents types de déclencheurs.

\subsubsection{Tableaux de bord avec Recharts}

Les tableaux de bord qualité utilisent \texttt{Recharts}~\cite{Recharts} pour la
visualisation des données~: graphiques en barres, courbes d'évolution, diagrammes
circulaires et jauges de performance.

%% ────────────────────────────────────────────────────────────────────────────
\subsection{Réalisation}
\label{subsec:sprint4_realisation}

\subsubsection{Interface d'évaluation qualité}

L'interface d'évaluation qualité permet au Quality Manager de visualiser la transcription,
le score IA et de saisir son évaluation manuelle.

\begin{figure}[H]
    \centering
    \fbox{\parbox{0.9\linewidth}{
    \centering\vspace{0.4cm}
    \textit{[Capture d'écran : Formulaire d'évaluation qualité]}\\[0.3cm]
    Cette capture montre le formulaire d'évaluation~: transcription de l'appel, score IA
    automatique affiché, champs de saisie pour la note manuelle, zone de commentaires
    et bouton de validation.
    \vspace{0.4cm}
    }}
    \caption{Formulaire d'évaluation qualité}
    \label{fig:sprint4_evaluation_form}
\end{figure}

\subsubsection{Dashboard qualité}

Le dashboard qualité offre une vue synthétique de la performance qualité du centre d'appels,
avec des graphiques et des indicateurs clés.

\begin{figure}[H]
    \centering
    \fbox{\parbox{0.9\linewidth}{
    \centering\vspace{0.4cm}
    \textit{[Capture d'écran : Dashboard qualité]}\\[0.3cm]
    Le dashboard qualité présente les métriques~: score qualité moyen, répartition des
    évaluations par agent, évolution temporelle des scores, alertes actives et
    comparaison scores IA vs manuels.
    \vspace{0.4cm}
    }}
    \caption{Dashboard qualité}
    \label{fig:sprint4_quality_dashboard}
\end{figure}

\subsubsection{Gestion des alertes}

L'interface de gestion des alertes permet de configurer les règles et de consulter
l'historique des alertes déclenchées.

\begin{figure}[H]
    \centering
    \fbox{\parbox{0.9\linewidth}{
    \centering\vspace{0.4cm}
    \textit{[Capture d'écran : Configuration des alertes]}\\[0.3cm]
    Interface de configuration des alertes avec les règles~: seuil de score minimum (70),
    délai d'inactivité (15 min), taux de conversion minimum (40\%). Chaque règle peut
    être activée ou désactivée.
    \vspace{0.4cm}
    }}
    \caption{Configuration des alertes}
    \label{fig:sprint4_alert_config}
\end{figure}

\subsubsection{Tests unitaires et validation}

Les tests du Sprint 4 couvrent~: le service de calcul de score qualité, l'évaluation
manuelle (CRUD), le système d'alertes (création, déclenchement, historique), les
tableaux de bord et les scénarios d'erreur.

\subsubsection{Résultats obtenus}

Le Sprint 4 complète le module qualité de la plateforme. Les Quality Managers disposent
d'outils pour évaluer manuellement les appels, configurer des alertes proactives et
suivre la performance qualité à travers des tableaux de bord interactifs. La combinaison
du scoring IA automatique et de l'évaluation manuelle offre une vision complète et
équilibrée de la qualité des interactions.

%% ══════════════════════════════════════════════════════════════════════════════
\section*{Conclusion}
\phantomsection
\label{sec:chap5_conclusion}

Ce chapitre a présenté le développement des fonctionnalités d'intelligence artificielle
et de gestion de la qualité de la plateforme CRM. Le Sprint 3 a permis de mettre en place
le coeur innovant du projet~: un pipeline IA complet combinant Whisper pour la transcription,
gemma3:4b pour le scoring et DistilBERT pour l'analyse de sentiment.

Le Sprint 4 a complété ce socle par des outils de gestion de la qualité~: évaluation manuelle,
alertes configurables et tableaux de bord interactifs. Ces fonctionnalités offrent aux
managers une vision complète de la performance qualité, combinant l'objectivité de l'IA
avec l'expertise humaine.

À l'issue de cette deuxième release, la plateforme dispose d'un système intégré d'analyse
et d'évaluation de la qualité des appels, positionnant la solution comme un CRM intelligent
de nouvelle génération. Les tests unitaires réalisés garantissent la fiabilité du pipeline
IA et des algorithmes de scoring.
